% Independent reconstruction of SECTIONS A.72; intended ID C.2.17.
\subsection{The prescribed finite quotient with the full abelian pair}

\begin{theorem}[Full reconstruction of SECTIONS A.72]
\label{thm:ghost-finite-quotient-survival}
Let \(N\) be an open subgroup of \(\Pi\), normal in
\(E_U\), let \(Q=\Pi/N\), and fix the actual section
\(s_3\) supplied by the rational point three.
There are an integer \(N_0>0\), disjoint sequences of
distinct primes \(q_j,r_j>3\), and an actual continuous
section
\[
 \tau:G_{\mathbb Q}\longrightarrow
 E_U/(N\cap\overline{[\Pi,\Pi]})
   =E_Q\mathbin{\times}_{E_{Q^{\rm ab}}}E_{\rm ab}
 \tag{FQ.1}
\]
with \(Q\)-component exactly \(s_3\) and abelian pair
\[
 A=3\prod_{j\ge0}q_j^{n_j},\qquad
 B=-2\prod_{j\ge0}r_j^{n_j},\qquad
 n_j=N_0(j+4)!.
 \tag{FQ.2}
\]
At every finite prime the pair is an ordinary
unit-equation pair, and the entire local section
in (FQ.1) is geometrically conjugate to the image
of that ordinary point. At the real place it has
the image of three. Both support sequences are
infinite. The proof specifies the permissible
parameters and the actual recursion, including
all its local component maps.
\end{theorem}
\begin{proof}
Write \(\Lambda=\Pi^{\rm ab}\), \(h=|Q|\), and
\(e=\exp(Q^{\rm ab})\).
The subgroup \(Ns_3(G_{\mathbb Q})\) defines a
finite etale cover \(f:Y\to U\) of degree \(h\),
geometrically the regular \(Q\)-cover, and an
actual rational point \(y_0\) above three.
The compactification is the normalization of
\(\mathbb P^1\) in \(\mathbb Q(Y)\).

\emph{Fixed fields and the exact intersection.}
Choose a finite Galois field \(F/\mathbb Q\)
containing the fields of definition of all
boundary points above zero and one, their local
leading coefficients, and the geometric deck
transformations. Include \(\mu_h\) and
\(3^{1/h},(-2)^{1/h}\).
At a zero-boundary point of ramification index
\(a\mid h\), choose a parameter with
\[
 t=d z^a(1+zh(z)),
 \tag{FQ.3}
\]
and include \(d^{1/a}\) in \(F\); use \(1-t\)
at a one-boundary point. There are finitely many
choices. Enlarge an integer \(\Delta\), divisible
by \(6h\), to give good integral models of the
cover, its boundary and these parameters outside
\(\Delta\).
At a prime outside \(\Delta\) split in \(F\),
the coordinate change
\(z\mapsto z(1+zh(z))^{1/a}\) is an invertible
power series on the maximal ideal: \(a\) is a
unit and the constant term of the root is one.
Thus every positive-valuation \(a\)-th power
value of \(t\) is obtained on this branch, with
all \(a\) root choices rational over the local
field. Applying this at every boundary point
proves that the entire \(Y\)-fibre splits for
such values; the corresponding assertion holds
at one.

Put \(M=[F:\mathbb Q]\), and form the fixed fields
\[
 K_d=F(\mu_{4M},d^{1/(4M)})\quad(d=3,-2).
 \tag{FQ.4}
\]
By Kronecker--Weber
\cite[I, Theorem 4.16]{GhostMilneCFT}, choose
conductors \(f_d\) with
\(K_d^{\rm ab}\subset\mathbb Q(\mu_{f_d})\).
Choose \(C\) divisible by \(24,4M,h\), and
\(N_0\) divisible by \(C,\Delta,f_3,f_{-2}\).
Every such choice is allowed below.

For \(24,4M,f_3,f_{-2}\mid n\), let
\(d'=d w^C\), where \(w>0\) is rational and
has valuation zero at two and three, and put
\(L=\mathbb Q(\mu_n,(d')^{1/n})\).
The radical calculation in the proof of
Theorem~\ref{thm:ghost-general-exponent-arithmetic}
applies to \(w^{C/24}\): the Kummer kernel
\(H=\operatorname{Gal}(L/\mathbb Q(\mu_n))\)
is cyclic of order \(n/2\).
Explicitly, the existing square root gives the
factor two, while an odd-prime root or a fourth
root would put a nonnormal rational pure-radical
field inside the abelian cyclotomic field.
Let \(H_F=\operatorname{Gal}(FL/F\mathbb Q(\mu_n))\).
Its index in \(H\) divides \(M\), whence
\(MH\subset H_F\). Complex conjugation acts by
inversion on \(H_F\), and therefore
\[
 2MH\subset2H_F\subset
 [\operatorname{Gal}(FL/\mathbb Q),
   \operatorname{Gal}(FL/\mathbb Q)].
 \tag{FQ.5}
\]
The fixed field of \(2MH\) in \(FL\) is
\[
 F\mathbb Q(\mu_n,(d')^{1/(4M)})
   =K_d\mathbb Q(\mu_n).
\]
For a Galois field \(K\) and an abelian field
\(T\), the Galois fibre product for \(KT\)
has commutator subgroup \([\operatorname{Gal}(K/\mathbb Q),
\operatorname{Gal}(K/\mathbb Q)]\times\{1\}\):
every element of the first factor lifts, and
commutators are computed coordinatewise.
Hence \((KT)^{\rm ab}=K^{\rm ab}T\).
Taking maximal abelian subfields in (FQ.5)
and using the conductors proves
\[
 (FL)^{\rm ab}=\mathbb Q(\mu_n).
 \tag{FQ.6}
\]
In particular, for every multiple \(P\) of
\(n\), \(FL\cap\mathbb Q(\mu_P)=\mathbb Q(\mu_n)\).
A unit CRT class modulo \(P\) that is one
modulo \(n\) agrees on this intersection with
the identity of \(FL\). The actual fibre
product and Chebotarev
\cite[Theorem 8.31]{GhostMilneANT} supply
arbitrarily large primes satisfying that CRT
class and splitting completely in \(FL\).

\emph{The designated component and a uniform local bound.}
For \(a=3u^C\), \(b=-2v^C\), with positive
rational \(u,v\), define
\[
 V_{a,b,n}:aX^n+bY^n=1,\quad XY\ne0,
 \qquad t=aX^n,
 \tag{FQ.7}
\]
where \(e\mid n\). The largest common geometric
quotient of the covers \(Y\) and \(V_{a,b,n}\)
is exactly \(Q^{\rm ab}\): any such quotient
is a quotient of \(\Lambda/n\Lambda\), hence
abelian, while \(e\mid n\) makes all of
\(Q^{\rm ab}\) available.
Let \(Z_Q=Y/[Q,Q]\), with its point-three
framing. The two coordinates
\[
 (a/3)^{1/e}X^{n/e},\qquad
 (b/(-2))^{1/e}Y^{n/e}
 \tag{FQ.8}
\]
give the rational morphism from \(V_{a,b,n}\)
to the pointed \(e\)-Kummer cover and then to
\(Z_Q\). Define the actual connected cover
\[
 W_{a,b,n}=Y\mathbin{\times}_{Z_Q}V_{a,b,n}.
 \tag{FQ.9}
\]
Its geometric group is
\(Q\times_{Q^{\rm ab}}\Lambda/n\Lambda\).
The common-quotient calculation proves geometric
connectedness. At \((a,b)=(3,-2)\) it contains
\((y_0,1,1)\).

Its degree \(d\) is at most \(D=hn^2\).
Riemann--Hurwitz, with at most three branch
values, gives \(2g\le d+2\), and the boundary
has at most \(3d\) geometric points.
At a prime \(p\nmid\Delta nab\), its
compactification has good reduction with the
same connected component and inertia data.
Indeed the Kummer covers by units extend
etale over the integral open curve; their
geometric order and that of \(Y\) are prime
to \(p\). Prime-to-\(p\) specialization for
the smooth pair \((\mathbb P^1,\{0,1,\infty\})\)
preserves these covers and inertia
\cite[Expos\'e XIII, Corollary 2.9,
\S2.10 and Lemma 2.11]{GhostSGA1}.
Locally at the boundary, normalization of the
tame equations in (FQ.3) gives smooth parameter
discs after an unramified extension, so the
compactification is smooth as asserted.
The trace formula and the weight-one eigenvalue
bound \cite[(1.5.1) and Theorem (1.6)]{GhostDeligneWeilI}
give at least
\(p+1-(D+2)\sqrt p-3D\) interior points.
This number is positive when
\[
 \begin{gathered}
 p>B_Q(n):=(4hn^2+10)^2,\\
 (4D+10)^2+1-(D+2)(4D+10)-3D
   =12D^2+59D+81>0.
 \end{gathered}
 \tag{FQ.10}
\]
Smooth Hensel lifting supplies a local point
of \(W_{a,b,n}\) whose \(X,Y\) are units.

\emph{The complete induction.}
Set \(n_j=N_0(j+4)!\), \((a_0,b_0)=(3,-2)\),
and \(W_j=W_{a_j,b_j,n_j}\).
We keep \(a_j=3u_j^C\), \(b_j=-2v_j^C\),
with positive integers \(u_j,v_j\) of disjoint
support away from six. Every \(W_j\) has a
unit-coordinate point at every finite prime.
At \(p\le B_Q(n_j)\) and at \(p\mid\Delta\)
we additionally keep a point over \(t=a_j\)
with \(X=1\), together with
\[
 \begin{gathered}
 Y\equiv1\pmod {p^{v_p(n_{j+1})}},\qquad
 V_j(p)=\max(v_p(a_j),v_p(b_j)),\\
 v_p(a_j+b_j-1)\ge
 V_j(p)+2v_p(n_{j+1})+j+3.
 \end{gathered}
 \tag{FQ.11}
\]
The rational point \((y_0,1,1)\) starts all
these conditions.

At stage \(j\), write \(n=n_j\), \(m=n_{j+1}\),
and let \(S\) contain the primes at most
\(B_Q(m)\), all divisors of \(\Delta\), and
all coefficient-support primes.
Choose an old point
\((y_p,\xi_p,\eta_p)\in W_j(\mathbb Q_p)\)
with unit coordinates for each \(p\in S\).
For \(p\mid m\) use (FQ.11).
Such a prime divides \(N_0\) or is at most
\(j+5\), and so is at most \(B_Q(n_j)\).
Consequently
\(\xi_p\equiv\eta_p\equiv1\pmod {p^{v_p(m)}}\).
Take precisions
\[
 L_p\ge V_j(p)+2v_p(n_{j+2})+j+4,
 \qquad P=\prod_{p\in S}p^{L_p}.
 \tag{FQ.12}
\]
They can be enlarged before selecting any prime
so that the following local continuation stays
in the designated component.

For unit variables \(Q_p,R_p\) close to
\(\xi_p,\eta_p\), the etale local inverse of
\(Y\to U\) gives \(y'_p\) above
\(a_jQ_p^n\), close to \(y_p\).
The polynomial
\[
 Z^m-\frac{1-a_jQ_p^n}{b_jR_p^n}
 \tag{FQ.13}
\]
has at one error valuation at least
\(L_p-v_p(b_j)>2v_p(m)\).
Hensel gives a root with
\(v_p(Z_p-1)\ge L_p-v_p(b_j)-v_p(m)\).
Put \(k=m/n\). Then
\((y'_p,Q_p,R_pZ_p^k)\) is a point of the
old fibre product, tending to
\((y_p,\xi_p,\eta_p)\).
Its designated component is open and closed;
increasing \(L_p\) therefore puts it in
\(W_j\). This uses etaleness in characteristic
zero even when \(p\mid e\); no finite-field
separation of those branches is needed.
The same estimate gives
\(Z_p\equiv1\pmod {p^{v_p(n_{j+2})}}\).

Choose a fresh prime \(q_j\), larger than all
members of \(S\) and than \(n_{j+2}\), with
\[
 q_j\equiv\xi_p\pmod {p^{L_p}}\ (p\in S),
 \qquad
 q_j\text{ split in }
 F\mathbb Q(\mu_n,b_j^{1/n}).
 \tag{FQ.14}
\]
The CRT class is one modulo \(m\), hence
modulo \(n\); (FQ.6) proves the required
compatibility. At \(q_j\) choose a unit
\(z\) with \(b_jz^n=1-a_jq_j^n\), by the
splitting and simple-root Hensel lifting.
Now choose a fresh prime \(r_j>q_j\) with
\[
 \begin{gathered}
 r_j\equiv\eta_p\pmod {p^{L_p}}\ (p\in S),
 \qquad r_j\equiv z\pmod {q_j},\\
 r_j\text{ split in }
 F\mathbb Q(\mu_m,(a_jq_j^n)^{1/m}).
 \end{gathered}
 \tag{FQ.15}
\]
Its combined CRT class is one modulo \(m\),
since \(q_j\nmid m\). Its radicand is still
three times a rational \(C\)-th power, because
\(C\mid n\). Thus (FQ.6) applies again.
Set \(a_{j+1}=a_jq_j^n\), \(b_{j+1}=b_jr_j^n\).
There is the literal morphism
\[
 W_{j+1}\longrightarrow W_j,\qquad
 (y,X,Y)\longmapsto
       (y,q_jX^{m/n},r_jY^{m/n}).
 \tag{FQ.16}
\]
The common coordinates (FQ.8) agree, since
\((a_{j+1}/3)^{1/e}=(a_j/3)^{1/e}q_j^{n/e}\),
and similarly for \(b\). The common quotient
remains \(Q^{\rm ab}\); hence its inverse image
is exactly the single geometrically connected
component (FQ.9) for the new coefficients.

For \(p\in S\), apply (FQ.13) with the actual
\(Q_p=q_j,R_p=r_j\). It gives the point
\((y'_p,1,Z_p)\) of \(W_{j+1}\).
The coefficient sum error has valuation at
least \(L_p\), and \(V_{j+1}(p)=V_j(p)\).
Thus (FQ.11) holds at the next stage.
At \(q_j\) set \(X=1\): (FQ.15) makes
\((1-a_jq_j^n)/(b_jr_j^n)\) congruent to one
modulo \(q_j\), so it has a unit \(m\)-th root.
The resulting value \(t=a_jq_j^n\) is an
\(a\)-th power of positive valuation for
every zero-boundary index \(a\mid h\), using
\(h\mid C,n\) and \(3^{1/h}\in F\).
By (FQ.3) its entire \(Y\)-fibre is rational.
Every common abelian sheet is available, so
one of its lifts is in the specified \(W_{j+1}\).
At \(r_j\) set \(Y=1\). The splitting in
(FQ.15) and Hensel give a unit \(X\), and
\(1-t=b_jr_j^n\) is an appropriate power at
every one-boundary branch. The same whole-fibre
argument selects a point of \(W_{j+1}\).
All remaining primes have unit coefficients,
do not divide \(\Delta m\), and exceed
\(B_Q(m)\); (FQ.10) supplies their points.
This completes every local case of the induction.

\emph{The actual limit and joint local conjugators.}
For each fixed finite modulus the factorial
exponents in (FQ.2), and also \(n_j/C\),
are eventually divisible by it. Their completed
products converge. Fix compatible roots of
each selected prime and use the literal relative
Kummer cocycles
\(\sum_j n_j\kappa(q_j)\),
\(\sum_j n_j\kappa(r_j)\).
Their tails vanish pointwise at every finite
coefficient level, so the sums are uniformly
convergent continuous cocycles with values in
\(C\Lambda\). Their image in \(Q^{\rm ab}\)
is zero. They consequently combine with the
actual \(Q\)-section \(s_3\) to give \(\tau\)
in (FQ.1).
For completeness, the group isomorphism in
(FQ.1) is the diagonal quotient map: its kernel
is the displayed intersection; surjectivity
follows by adjusting two representatives whose
difference lies in
\(N\overline{[\Pi,\Pi]}\). Normality and
compactness make this product a closed subgroup.

The finite geometric quotient at exponent
\(n_j\) is
\(Q\times_{Q^{\rm ab}}\Lambda/n_j\Lambda\),
and the cover defined by \(\tau\) there is
exactly \(W_j\): its abelian twist is
\((a_j,b_j)\), and its common coordinates are
precisely (FQ.8). The transition (FQ.16)
adds \(n_j\kappa(q_j),n_j\kappa(r_j)\),
which vanish at that exponent. Thus the
quotient sections and their framings agree.

For a fixed finite prime \(p\), the ordinary
products of the local units converge: \(n_j\)
eventually annihilates every finite quotient
of \(\mathbb Z_p^\times\), including at two.
Their valuations stabilize, so the limits
\(a_j\to x_p\), \(b_j\to y_p\) are nonzero.
Equation (FQ.11) gives \(x_p+y_p=1\).
These are the localizations of the completed
products, by the same finite-power comparisons.
Choose any unit-coordinate local point on
each \(W_j\), with image \(t_j=a_jX_j^{n_j}\).
Uniform annihilation of finite unit quotients
gives \(t_j\to x_p\).
For fixed \(i\), (FQ.16) supplies a lift of
every \(t_j\), \(j\ge i\), to \(W_i\).
The image of \(W_i(\mathbb Q_p)\) is closed
in \(U(\mathbb Q_p)\): pass to its projective
compactification, whose local points are
compact, and use that \(x_p\ne0,1,\infty\).
Hence \(x_p\) lifts to every \(W_i\).

A fixed sheet above \(x_p\) is exactly a
geometric conjugacy of its point section with
\(\tau\) in that finite quotient. Pull these
conjugator sets back to the compact geometric
group in (FQ.1). They are nonempty, closed
and nested by (FQ.16); their intersection
gives one conjugator for the whole section,
since the exponents \(n_i\) are cofinal.
At the real place choose positive compatible
roots for all the divided products. Their
relative cocycles vanish on complex conjugation,
so the real section is the image of three.
Finally each new prime is distinct and retains
its displayed nonzero integer valuation.
The infinite supports in (FQ.2) follow, and
finite-support reconstruction excludes a
rational point image. This proves every
assertion of the theorem.
\end{proof}

The independent certificate
\texttt{check\_finite\_quotient\_survival.py}
checks \(2520\) cyclic degree/index cases in
(FQ.5), the exact polynomial bound (FQ.10),
and \(546\) common-coordinate exponent transitions.
It also computes actual roots of (FQ.13) to
the stated precision at \(p=2,3,5,7,11,13\),
including the nonunit derivatives at two,
three and five. These calculations verify
the finite arithmetic steps; the field and
cover existence statements are proved above.
